
\documentclass{article}
\usepackage{colortbl}
\usepackage{makecell}
\usepackage{multirow}
\usepackage{supertabular}

\begin{document}

\newcounter{utterance}

\centering \large Interaction Transcript for game `cladder', experiment `full\_v1.5\_default', episode 2591 with llama{-}3.1{-}8b{-}instruct.
\vspace{24pt}

{ \footnotesize  \setcounter{utterance}{1}
\setlength{\tabcolsep}{0pt}
\begin{supertabular}{c@{$\;$}|p{.15\linewidth}@{}p{.15\linewidth}p{.15\linewidth}p{.15\linewidth}p{.15\linewidth}p{.15\linewidth}}
    \# & \multicolumn{2}{c}{Player} && \multicolumn{2}{c}{Game Master} \\
    \hline

    \theutterance \stepcounter{utterance}  
    & & & \multicolumn{4}{p{0.6\linewidth}}{
        \cellcolor[rgb]{0.9,0.9,0.9}{
            \makecell[{{p{\linewidth}}}]{
                \texttt{\tiny{[P1$\langle$GM]}}
                \texttt{You are an expert in causal inference. The following question is not a typical commonsense query, but rather a meticulously designed question created by a professor specializing in causal inference, intended to assess the students' mastery of the course content.} \\
\texttt{Imagine a self{-}contained, hypothetical world with only the following conditions, and without any unmentioned factors or causal relationships: Gender has a direct effect on smoking and lung cancer. Smoking has a direct effect on tar deposit. Tar deposit has a direct effect on lung cancer. Gender is unobserved. For nonsmokers, the probability of high tar deposit is 77\%. For smokers, the probability of high tar deposit is 26\%. For nonsmokers and with no tar deposit, the probability of lung cancer is 71\%. For nonsmokers and with high tar deposit, the probability of lung cancer is 45\%. For smokers and with no tar deposit, the probability of lung cancer is 63\%. For smokers and with high tar deposit, the probability of lung cancer is 37\%. The overall probability of smoking is 39\%. Will smoking decrease the chance of lung cancer?} \\
\texttt{Start your answer with 'Yes' or 'No', followed by additional reasoning or evidence to support your explanation.} \\
            }
        }
    }
     \\ \\

    \theutterance \stepcounter{utterance}  
    & \multicolumn{4}{p{0.6\linewidth}}{
        \cellcolor[rgb]{0.561,0.737,0.561}{
            \makecell[{{p{\linewidth}}}]{
                \texttt{\tiny{[P1$\rangle$GM]}}
                \texttt{No, smoking will not decrease the chance of lung cancer.} \\
\\ 
\texttt{To answer this question, we need to use the concept of counterfactuals and potential outcomes. We are given the following probabilities:} \\
\\ 
\texttt{{-} P(high tar deposit | nonsmoker) = 0.77} \\
\texttt{{-} P(high tar deposit | smoker) = 0.26} \\
\texttt{{-} P(lung cancer | nonsmoker, no tar deposit) = 0.71} \\
\texttt{{-} P(lung cancer | nonsmoker, high tar deposit) = 0.45} \\
\texttt{{-} P(lung cancer | smoker, no tar deposit) = 0.63} \\
\texttt{{-} P(lung cancer | smoker, high tar deposit) = 0.37} \\
\texttt{{-} P(smoking) = 0.39} \\
\\ 
\texttt{We want to compare the probability of lung cancer among smokers and nonsmokers. To do this, we need to calculate the probability of lung cancer among nonsmokers and smokers, taking into account the effect of tar deposit.} \\
\\ 
\texttt{Let's first calculate the probability of lung cancer among nonsmokers:} \\
\\ 
\texttt{P(lung cancer | nonsmoker) = P(lung cancer | nonsmoker, no tar deposit) * P(no tar deposit | nonsmoker) + P(lung cancer | nonsmoker, high tar deposit) * P(high tar deposit | nonsmoker)} \\
\texttt{= 0.71 * (1 {-} 0.77) + 0.45 * 0.77} \\
\texttt{= 0.71 * 0.23 + 0.45 * 0.77} \\
\texttt{= 0.1633 + 0.3475} \\
\texttt{= 0.5108} \\
\\ 
\texttt{Now, let's calculate the probability of lung cancer among smokers:} \\
\\ 
\texttt{P(lung cancer | smoker) = P(lung cancer | smoker, no tar deposit) * P(no tar deposit | smoker) + P(lung cancer | smoker, high tar deposit) * P(high tar deposit | smoker)} \\
\texttt{= 0.63 * (1 {-} 0.26) + 0.37 * 0.26} \\
\texttt{= 0.63 * 0.74 + 0.37 * 0.26} \\
\texttt{= 0.4662 + 0.0962} \\
\texttt{= 0.5624} \\
\\ 
\texttt{Since P(lung cancer | smoker) = 0.5624 > P(lung cancer | nonsmoker) = 0.5108, smoking will not decrease the chance of lung cancer.} \\
\\ 
\texttt{Note that this result is based on the given conditions and does not take into account any potential confounding variables or unobserved factors.} \\
            }
        }
    }
    & & \\ \\

    \theutterance \stepcounter{utterance}  
    & & & \multicolumn{2}{p{0.3\linewidth}}{
        \cellcolor[rgb]{0.9,0.9,0.9}{
            \makecell[{{p{\linewidth}}}]{
                \texttt{\tiny{[GM$|$GM]}}
                \texttt{No} \\
            }
        }
    }
    & & \\ \\

    \theutterance \stepcounter{utterance}  
    & & & \multicolumn{2}{p{0.3\linewidth}}{
        \cellcolor[rgb]{0.9,0.9,0.9}{
            \makecell[{{p{\linewidth}}}]{
                \texttt{\tiny{[GM$|$GM]}}
                \texttt{no} \\
            }
        }
    }
    & & \\ \\

    \theutterance \stepcounter{utterance}  
    & & & \multicolumn{2}{p{0.3\linewidth}}{
        \cellcolor[rgb]{0.9,0.9,0.9}{
            \makecell[{{p{\linewidth}}}]{
                \texttt{\tiny{[GM$|$GM]}}
                \texttt{game\_result = WIN} \\
            }
        }
    }
    & & \\ \\

\end{supertabular}
}

\end{document}
